% GENERATED FILE. Edit canonical lessons, then run npm run generate:books.
% canonical-source-hash: fnv1a64:726d7606ea1fe184
% canonical-lessons: SA-C487-dantadhavanam, SA-C487-prasadhanam, SA-C487-anjanam, SA-C487-angaragah, SA-C487-paridadhati, SA-R487-review-480-487

\chapter{Getting Ready}
\label{ch:sa-getting-ready}
\hlchaptermodality{487}

\begin{chapteropening}
\textbf{By the end of this chapter:} \emph{I can talk about grooming and dressing up.}

Complete the last lesson of chapter 487: I can talk about grooming and dressing up.
\end{chapteropening}

\section[\texorpdfstring{\sk{दन्तधावनम्} (dantadhāvanam)}{दन्तधावनम् (dantadhāvanam)}]{\sk{दन्तधावनम्} (dantadhāvanam) — brushing the teeth}

\label{lesson:SA-C487-dantadhavanam}



\begin{quote}
Before the new one: say the Sanskrit for an oil massage, then the Sanskrit for the hour before dawn.
\end{quote}

\subsection*{You'll want to know: \sk{दन्तधावनम्}}
\textbf{\sk{दन्तधावनम्}} — \emph{dantadhāvanam} — \textquotedblleft{}brushing the teeth\textquotedblright{}.

\begin{quote}
\sk{प्रातः} \sk{दन्तधावनं} \sk{करोतु।}
\end{quote}

\begin{quote}
\emph{prātaḥ dantadhāvanaṃ karotu.}
\end{quote}

\begin{quote}
Brush your teeth in the morning.
\end{quote}

A neuter noun.

\begin{cousinweb}
\textbf{\sk{दन्तधावनम्}} (\emph{dantadhāvanam}) is built from \textbf{\sk{दन्तः}}, \textquotedblleft{}a tooth\textquotedblright{}, and \emph{dhāvana}, \textquotedblleft{}cleansing\textquotedblright{}, from the root √dhāv, \textquotedblleft{}to rinse\textquotedblright{}; MW records it from the Kauśika Sūtra on (MW).
\end{cousinweb}

\begin{grammarlens}[title={What you've built}]
One more word for describing what you lived and felt.
\end{grammarlens}

\subsection*{Your turn}
\begin{itemize}
\raggedright
  \item \emph{Say it:} \emph{dantadhāvanam}
  \item \emph{Say it:} \emph{dantadhāvanam}, once more
  \item \emph{Say it:} \emph{prātaḥ dantadhāvanaṃ karotu.} — Brush your teeth in the morning.
  \item \emph{Recall:} say the Sanskrit for an oil massage, then the Sanskrit for the hour before dawn, then say \emph{dantadhāvanam} again
\end{itemize}

\begin{tcolorbox}[breakable,colback=teal!4,colframe=teal!35!black,title={Before you move on}]
What does \sk{दन्तधावनम्} mean? (\textquotedblleft{}Brushing the teeth\textquotedblright{}.) Say it once more.
\end{tcolorbox}
\section[\texorpdfstring{\sk{प्रसाधनम्} (prasādhanam)}{प्रसाधनम् (prasādhanam)}]{\sk{प्रसाधनम्} (prasādhanam) — decoration; dressing up}

\label{lesson:SA-C487-prasadhanam}



\begin{quote}
Before the new one: say the Sanskrit for the hour before dawn, then the Sanskrit for brushing the teeth.
\end{quote}

\subsection*{You'll want to know: \sk{प्रसाधनम्}}
\textbf{\sk{प्रसाधनम्}} — \emph{prasādhanam} — \textquotedblleft{}decoration; dressing up\textquotedblright{}.

\begin{quote}
\sk{उत्सवात्} \sk{पूर्वं} \sk{मीरा} \sk{प्रसाधनं} \sk{करोति।}
\end{quote}

\begin{quote}
\emph{utsavāt pūrvaṃ mīrā prasādhanaṃ karoti.}
\end{quote}

\begin{quote}
Before the festival Mira dresses up.
\end{quote}

A neuter noun.

\begin{cousinweb}
\textbf{\sk{प्रसाधनम्}} (\emph{prasādhanam}) is built from \emph{pra} and the root √sādh, \textquotedblleft{}to accomplish\textquotedblright{}: a making ready; MW gives \textquotedblleft{}embellishment, decoration, toilet\textquotedblright{} (MW).
\end{cousinweb}

\begin{grammarlens}[title={What you've built}]
One more word for describing what you lived and felt.
\end{grammarlens}

\subsection*{Your turn}
\begin{itemize}
\raggedright
  \item \emph{Say it:} \emph{prasādhanam}
  \item \emph{Say it:} \emph{prasādhanam}, once more
  \item \emph{Say it:} \emph{utsavāt pūrvaṃ mīrā prasādhanaṃ karoti.} — Before the festival Mira dresses up.
  \item \emph{Recall:} say the Sanskrit for the hour before dawn, then the Sanskrit for brushing the teeth, then say \emph{prasādhanam} again
\end{itemize}

\begin{tcolorbox}[breakable,colback=teal!4,colframe=teal!35!black,title={Before you move on}]
What does \sk{प्रसाधनम्} mean? (\textquotedblleft{}Decoration; dressing up\textquotedblright{}.) Say it once more.
\end{tcolorbox}
\section[\texorpdfstring{\sk{अञ्जनम्} (añjanam)}{अञ्जनम् (añjanam)}]{\sk{अञ्जनम्} (añjanam) — kohl}

\label{lesson:SA-C487-anjanam}



\begin{quote}
Before the new one: say the Sanskrit for brushing the teeth, then the Sanskrit for decoration.
\end{quote}

\subsection*{You'll want to know: \sk{अञ्जनम्}}
\textbf{\sk{अञ्जनम्}} — \emph{añjanam} — \textquotedblleft{}kohl\textquotedblright{}.

\begin{quote}
\sk{नारी} \sk{नेत्रे} \sk{अञ्जनं} \sk{स्थापयति।}
\end{quote}

\begin{quote}
\emph{nārī netre añjanaṃ sthāpayati.}
\end{quote}

\begin{quote}
The woman puts kohl on her eyes.
\end{quote}

A neuter noun.

\begin{cousinweb}
\textbf{\sk{अञ्जनम्}} (\emph{añjanam}) comes from the root √añj, \textquotedblleft{}to anoint\textquotedblright{}; MW gives \textquotedblleft{}black pigment or collyrium applied to the eyelashes\textquotedblright{} (MW).
\end{cousinweb}

\begin{grammarlens}[title={What you've built}]
One more word for describing what you lived and felt.
\end{grammarlens}

\subsection*{Your turn}
\begin{itemize}
\raggedright
  \item \emph{Say it:} \emph{añjanam}
  \item \emph{Say it:} \emph{añjanam}, once more
  \item \emph{Say it:} \emph{nārī netre añjanaṃ sthāpayati.} — The woman puts kohl on her eyes.
  \item \emph{Recall:} say the Sanskrit for brushing the teeth, then the Sanskrit for decoration, then say \emph{añjanam} again
\end{itemize}

\begin{tcolorbox}[breakable,colback=teal!4,colframe=teal!35!black,title={Before you move on}]
What does \sk{अञ्जनम्} mean? (\textquotedblleft{}Kohl\textquotedblright{}.) Say it once more.
\end{tcolorbox}
\section[\texorpdfstring{\sk{अंगरागः} (aṅgarāgaḥ)}{अंगरागः (aṅgarāgaḥ)}]{\sk{अंगरागः} (aṅgarāgaḥ) — a scented cosmetic}

\label{lesson:SA-C487-angaragah}



\begin{quote}
Before the new one: say the Sanskrit for decoration, then the Sanskrit for kohl.
\end{quote}

\subsection*{You'll want to know: \sk{अंगरागः}}
\textbf{\sk{अंगरागः}} — \emph{aṅgarāgaḥ} — \textquotedblleft{}a scented cosmetic\textquotedblright{}.

\begin{quote}
\sk{स्नानात्} \sk{अनन्तरं} \sk{सा} \sk{अंगरागं} \sk{लिम्पति।}
\end{quote}

\begin{quote}
\emph{snānāt anantaraṃ sā aṅgarāgaṃ limpati.}
\end{quote}

\begin{quote}
After her bath she puts on a scented cream.
\end{quote}

A masculine noun.

\begin{cousinweb}
\textbf{\sk{अंगरागः}} (\emph{aṅgarāgaḥ}) is built from \textbf{\sk{अंगम्}}, \textquotedblleft{}a limb, the body\textquotedblright{}, and \emph{rāga}, \textquotedblleft{}colour\textquotedblright{}; MW gives \textquotedblleft{}application of unguents or cosmetics to the body\textquotedblright{} and \textquotedblleft{}scented cosmetic\textquotedblright{} (MW).
\end{cousinweb}

\begin{grammarlens}[title={What you've built}]
One more word for describing what you lived and felt.
\end{grammarlens}

\subsection*{Your turn}
\begin{itemize}
\raggedright
  \item \emph{Say it:} \emph{aṅgarāgaḥ}
  \item \emph{Say it:} \emph{aṅgarāgaḥ}, once more
  \item \emph{Say it:} \emph{snānāt anantaraṃ sā aṅgarāgaṃ limpati.} — After her bath she puts on a scented cream.
  \item \emph{Recall:} say the Sanskrit for decoration, then the Sanskrit for kohl, then say \emph{aṅgarāgaḥ} again
\end{itemize}

\begin{tcolorbox}[breakable,colback=teal!4,colframe=teal!35!black,title={Before you move on}]
What does \sk{अंगरागः} mean? (\textquotedblleft{}A scented cosmetic\textquotedblright{}.) Say it once more.
\end{tcolorbox}
\section[\texorpdfstring{\sk{परिदधाति} (paridadhāti)}{परिदधाति (paridadhāti)}]{\sk{परिदधाति} (paridadhāti) — puts on clothes; wears}

\label{lesson:SA-C487-paridadhati}



\begin{quote}
Before the new one: say the Sanskrit for kohl, then the Sanskrit for a scented cosmetic.
\end{quote}

\subsection*{You'll want to know: \sk{परिदधाति}}
\textbf{\sk{परिदधाति}} — \emph{paridadhāti} — \textquotedblleft{}puts on clothes; wears\textquotedblright{}.

\begin{quote}
\sk{उत्सवे} \sk{मीरा} \sk{नवं} \sk{वस्त्रं} \sk{परिदधाति।}
\end{quote}

\begin{quote}
\emph{utsave mīrā navaṃ vastraṃ paridadhāti.}
\end{quote}

\begin{quote}
At the festival Mira puts on new clothes.
\end{quote}

Past (she …): \sk{सा} \sk{नवं} \sk{वस्त्रं} \sk{परिहितवती}, \emph{sā navaṃ vastraṃ parihitavatī}.

\begin{cousinweb}
\textbf{\sk{परिदधाति}} (\emph{paridadhāti}) is built from \emph{pari}, \textquotedblleft{}around\textquotedblright{}, and the root √dhā, \textquotedblleft{}to put\textquotedblright{}: to put round oneself; MW gives \textquotedblleft{}to put on, wear, dress\textquotedblright{}, and its ppp \emph{parihita} is \textquotedblleft{}put on, clothed\textquotedblright{} (MW).
\end{cousinweb}

\begin{grammarlens}[title={What you've built}]
One more word for describing what you lived and felt.
\end{grammarlens}

\subsection*{Your turn}
\begin{itemize}
\raggedright
  \item \emph{Say it:} \emph{paridadhāti}
  \item \emph{Say it:} \emph{paridadhāti}, once more
  \item \emph{Say it:} \emph{utsave mīrā navaṃ vastraṃ paridadhāti.} — At the festival Mira puts on new clothes.
  \item \emph{Recall:} say the Sanskrit for kohl, then the Sanskrit for a scented cosmetic, then say \emph{paridadhāti} again
\end{itemize}

\begin{tcolorbox}[breakable,colback=teal!4,colframe=teal!35!black,title={Before you move on}]
What does \sk{परिदधाति} mean? (\textquotedblleft{}Puts on clothes; wears\textquotedblright{}.) Say it once more.
\end{tcolorbox}
\section[(dialogue)]{Looking back: from Bites and Wounds to Getting Ready}

\label{lesson:SA-R487-review-480-487}



\begin{quote}
Say the Sanskrit for bites and poison.
\end{quote}

\subsection*{Your turn}
Say these aloud:

\begin{itemize}
\raggedright
  \item bites, poison, an injury, a broken bone, oozes
  \item blind, deaf, mute, lame, hunchbacked
  \item treats a patient, a patient, the cause of an illness, medical treatment, a splinter
  \item a pill, a powder, an ointment, a decoction, smears
  \item the science of health, the wind humour, bile, phlegm, a wholesome diet
  \item nursing, attends on, revives, relief, a rejuvenating tonic
  \item fasts, breath control, meditation, an oil massage, the hour before dawn
  \item brushing the teeth, decoration, kohl, a scented cosmetic, puts on clothes
\end{itemize}

\begin{tcolorbox}[breakable,colback=teal!4,colframe=teal!35!black,title={Before you move on}]
Bites? (\textbf{\sk{दशति}}, \emph{daśati}.) Poison? (\textbf{\sk{विषम्}}, \emph{viṣam}.) An injury? (\textbf{\sk{क्षतम्}}, \emph{kṣatam}.) A broken bone? (\textbf{\sk{अस्थिभंगः}}, \emph{asthibhaṅgaḥ}.) Oozes? (\textbf{\sk{स्रवति}}, \emph{sravati}.) Blind? (\textbf{\sk{अन्धः}}, \emph{andhaḥ}.) Deaf? (\textbf{\sk{बधिरः}}, \emph{badhiraḥ}.) Mute? (\textbf{\sk{मूकः}}, \emph{mūkaḥ}.) Lame? (\textbf{\sk{पंगुः}}, \emph{paṅguḥ}.) Hunchbacked? (\textbf{\sk{कुब्जः}}, \emph{kubjaḥ}.) Treats a patient? (\textbf{\sk{चिकित्सति}}, \emph{cikitsati}.) A patient? (\textbf{\sk{आतुरः}}, \emph{āturaḥ}.) The cause of an illness? (\textbf{\sk{निदानम्}}, \emph{nidānam}.) Medical treatment? (\textbf{\sk{उपचारः}}, \emph{upacāraḥ}.) A splinter? (\textbf{\sk{शल्यम्}}, \emph{śalyam}.) A pill? (\textbf{\sk{गुटिका}}, \emph{guṭikā}.) A powder? (\textbf{\sk{चूर्णम्}}, \emph{cūrṇam}.) An ointment? (\textbf{\sk{लेपः}}, \emph{lepaḥ}.) A decoction? (\textbf{\sk{कषायः}}, \emph{kaṣāyaḥ}.) Smears? (\textbf{\sk{लिम्पति}}, \emph{limpati}.) The science of health? (\textbf{\sk{आयुर्वेदः}}, \emph{āyurvedaḥ}.) The wind humour? (\textbf{\sk{वातः}}, \emph{vātaḥ}.) Bile? (\textbf{\sk{पित्तम्}}, \emph{pittam}.) Phlegm? (\textbf{\sk{कफः}}, \emph{kaphaḥ}.) A wholesome diet? (\textbf{\sk{पथ्यम्}}, \emph{pathyam}.) Nursing? (\textbf{\sk{शुश्रूषा}}, \emph{śuśrūṣā}.) Attends on? (\textbf{\sk{उपचरति}}, \emph{upacarati}.) Revives? (\textbf{\sk{आप्याययति}}, \emph{āpyāyayati}.) Relief? (\textbf{\sk{उपशमः}}, \emph{upaśamaḥ}.) A rejuvenating tonic? (\textbf{\sk{रसायनम्}}, \emph{rasāyanam}.) Fasts? (\textbf{\sk{उपवसति}}, \emph{upavasati}.) Breath control? (\textbf{\sk{प्राणायामः}}, \emph{prāṇāyāmaḥ}.) Meditation? (\textbf{\sk{ध्यानम्}}, \emph{dhyānam}.) An oil massage? (\textbf{\sk{अभ्यंगः}}, \emph{abhyaṅgaḥ}.) The hour before dawn? (\textbf{\sk{ब्रह्ममुहूर्तः}}, \emph{brahmamuhūrtaḥ}.) Brushing the teeth? (\textbf{\sk{दन्तधावनम्}}, \emph{dantadhāvanam}.) Decoration? (\textbf{\sk{प्रसाधनम्}}, \emph{prasādhanam}.) Kohl? (\textbf{\sk{अञ्जनम्}}, \emph{añjanam}.) A scented cosmetic? (\textbf{\sk{अंगरागः}}, \emph{aṅgarāgaḥ}.) Puts on clothes? (\textbf{\sk{परिदधाति}}, \emph{paridadhāti}.)
\end{tcolorbox}
